
\chapter*{Acknowledgments}
I would like to express my deepest gratitude to my thesis supervisor, Águeda Gómez Cambronero. Thank you for proposing this project and for your invaluable guidance, continuous support, and technical advice throughout the development of \textit{Micorrizas}, but above all, thank you for always being available and for all the time you have invested in me. Your help with the implementation and the resources you provided have been fundamental to the success of this work.

I am also immensely grateful to Silvia Ferrer Castillo and Marta Mestre Aguila for their essential contributions to the project’s pedagogical and educational framework. Both their initial proposal and their subsequent guidance were key to shaping the educational value of the experience.

Thanks to Fabio Nicolas Luna Arevalo for his technical assistance, specifically for providing access and support with the ArcGIS SDK, which made it possible to create a very accurate map of the \textit{Jardín de los Sentidos}.

I would also like to thank Luis De Quevedo for create the game logo, and Carlos Alonso for creating the gameplay video. Their contributions helped to improve the presentation of the project and to communicate the final result more clearly.

Finally, I am also grateful to the friends who took part in the playtesting sessions. Their time, involvement and feedback were very valuable for detecting issues, refining the experience and improving the final version of the game.

\clearpage

\chapter*{Abstract}
This document presents the Bachelor’s Thesis developed for the Bachelor’s Degree in Video Game Design and Development. The project originates from a pedagogical activity carried out with Early Childhood Education students in the course MI1815 – The Natural Environment in Early Childhood Education, consisting of a guided visit to the \textit{Jardín de los Sentidos}, an outdoor space at Jaume I University characterised by its rich biodiversity. The educational value of this activity lies in encouraging students to observe, discuss, and reflect on natural ecosystems through direct interaction with the environment.
Building upon this experience, this project explores how educational objectives of the original activity can be transferred to a digital, autonomous and playful format. To address this challenge, \textit{Micorrizas} was designed as an educational video game that promotes environmental awareness and engagement through active exploration and collaborative learning.
The resulting game is a cooperative location-based mobile experience developed with Unity, where players interact with real-world locations while progressing through shared challenges and decision-making processes. To support this experience, the game combines geolocation and multiplayer technologies that enable collaborative gameplay in outdoor environments. The game has also been evaluated by the lecturers responsible for the original educational activity, providing feedback on its suitability as a tool for environmental education in outdoor learning contexts.

\vspace{1cm}
\section*{Keywords}
Serious game - educational game - cooperative multiplayer - geolocation - mobile learning - environmental education - Jardín de los Sentidos.

\clearpage

\tableofcontents
\clearpage

\listoffigures
\clearpage

\listoftables
\clearpage

